\section{Theory}
\label{sec:Theorie}
The theoretical backgrounds for this experiment are discussed in this section. 


\subsection{Rutherford scattering formula}
First the Rutherford scattering formula is needed. 
If alpha particles are scattered at a foil with thickness $D$ of the material with the nuclear charge number $Z$, the number of scattered particles $\mathrm{d} n$ depending 
on the scattering angle $\theta$ in the solid angle $\mathrm{d}\Omega$ is given by
\begin{equation}
    \frac{\mathrm{d} n}{n} = \frac{Z^2 e^2 D N}{\left(2 \pi \epsilon_0\right)^2 E_{\mathrm{kin}}^2 \mathrm{sin}\left(\frac{\theta}{2}\right)} \, \mathrm{d} \Omega \, .
    \label{eqn:rutherford}
\end{equation}
Here $n$ is the total number of scattered particles, $e$ is the elementary charge, $epsilon_0$ is the permittivity constant of vacuum and $E_{\mathrm{kin}}= \frac{1}{2} m v_o^2$ is the 
kinetic energy of the alpha particles \cite[39-44]{haken_1984}.

\subsection{Decay chain of the used radioactive material}
In this experiment americium is used as a radioactive source of alpha particles. 


\subsection{Surface barrier detector}
\cite{leo_1994}

\subsection{Uncertainties of the measurement}
\label{sec:fehlerrechnung}
To handle the uncertainties of the measurements the folowing formulas are used:
The Number of measured events is assumed to follow the poisson distribution. 
Therefor the uncertainty for $N$ measured events is given by $\mathrm{\Delta} N = \sqrt{N}$
It follows that the relativ uncertainty can be calculatet using %QUELLEEEEEEEEEEEEE!!!!!!!!!
\begin{equation}
    \frac{\mathrm{\Delta}N}{N} = \frac{1}{\sqrt{N}} \, .
    \label{eqn:relative_poisson}
\end{equation}
\\
%\subsection{Mittelwert}
The mean value is calculatet with
\begin{equation}
\bar{x} = \frac{1}{N} \sum_{i=1}^{N} x_i \, .
\label{eqn:mittelwert}
\end{equation}
%\subsection{Standardabweichung}
\\
The standard deviation is given by
\begin{equation}
\sigma = \sqrt{\frac{1}{N - 1} \sum_{i=1}^{N} \left( x_i - \bar{x}_N \right)^2} \, .
\label{eqn:standardabweichung}
\end{equation}
berechnet.
%%\subsubsection{Standardabweichung des Mittelwerts}
\\
The standard deviation of the mean value is given by
\begin{equation}
\Delta \bar{x} = \frac{\sigma}{\sqrt{N}} \, .
\label{eqn:standardmittelwert}
\end{equation}
%%\subsection{Gaußsche Fehlerfortpflanzung}
\\
The deveation of a value derived from experimental parameters can be calculatet using the Gaussian error propagation
\begin{equation}
\Delta f = \sqrt{\sum_{i=1}^{N} \left( \frac{\partial f}{\partial x_i} \right)^2 \cdot (\Delta x_i)^2} \, .
\label{eqn:fehlerfortpflanzung}
\end{equation}
\\
The relativ deviation of a experimental parameter from the theoretical parameter is given by ~\cite{relabw}
\begin{equation}
    \increment x_{\mathrm{rel}} = \left| \frac{x_{\mathrm{\exp}}-x_{\mathrm{theo}}}{x_{\mathrm{theo}}}\right|\, .
    \label{eqn:relabw}
\end{equation}

